%%%%%%%%%%%%%%%%%%%%%%%%%%%%%%%%%%%%%%%%%%%%%%%%%%%%%%%%%%%%%%%%%%%
%
% New Commands
%
%%%%%%%%%%%%%%%%%%%%%%%%%%%%%%%%%%%%%%%%%%%%%%%%%%%%%%%%%%%%%%%%%%%
%------------------------------------------------------------------
\newcommand{\overbar}[1]{\mkern 1.5mu\overline{\mkern-1.5mu#1\mkern-1.5mu}\mkern 1.5mu}
\newcommand{\unirand}{\ensuremath{\stackrel{\mathrm{\$}}{\leftarrow}}} %Uniform Random
\newcommand{\textunderscript}[1]{$_{\text{#1}}$}
%------------------------------------------------------------------

%%%%%%%%%%%%%%%%%%%%%%%%%%%%%%%%%%%%%%%%%%%%%%%%%%%%%%%%%%%%%%%%%%%
%
% Define Math Symbol
%
%%%%%%%%%%%%%%%%%%%%%%%%%%%%%%%%%%%%%%%%%%%%%%%%%%%%%%%%%%%%%%%%%%%
%------------------------------------------------------------------
\DeclarePairedDelimiter{\ceil}{\lceil}{\rceil}
\captionsetup{compatibility=false}
%------------------------------------------------------------------

%%%%%%%%%%%%%%%%%%%%%%%%%%%%%%%%%%%%%%%%%%%%%%%%%%%%%%%%%%%%%%%%%%%
%
% Define Theorem
%
%%%%%%%%%%%%%%%%%%%%%%%%%%%%%%%%%%%%%%%%%%%%%%%%%%%%%%%%%%%%%%%%%%%
%------------------------------------------------------------------
\long\def\comment#1{}
\newtheorem{theorem}{Theorem}
\newtheorem{proposition}{Proposition}
\newtheorem{corollary}{Corollary}
\newtheorem{lemma}{Lemma}
\newtheorem{definition}{Definition}
\def\qed{\hfill $\Box$}
%------------------------------------------------------------------

%%%%%%%%%%%%%%%%%%%%%%%%%%%%%%%%%%%%%%%%%%%%%%%%%%%%%%%%%%%%%%%%%%%
%
% Comment
%
%%%%%%%%%%%%%%%%%%%%%%%%%%%%%%%%%%%%%%%%%%%%%%%%%%%%%%%%%%%%%%%%%%%
%------------------------------------------------------------------
% \renewcommand{\baselinestretch}{2}
\newcommand{\ksedit}[1]{\textcolor{blue}{#1}}
\renewcommand{\ksedit}[1]{#1}
%\renewcommand{\ksedit}[1]{}
\newcommand{\sgnote}[1]{\textcolor{red}{\textbf{[SG: #1]}}}
\renewcommand{\sgnote}[1]{}
%------------------------------------------------------------------

%% Note for editor/advisor, e.g. \ednote{Add a new paragraph}
\newcommand{\ednote}[1]{\textbf{Edit: #1}}

%%% For writing the vector
\newcommand*{\vv}[1]{\vec{\mkern0mu#1}}

%%% Somsak's macros
%% For biber. Some of the commands in Bibtex such as \citep and \cite are not defined
%% or behave differently in Biber. 
\newcommand{\citep}{\parencite}
\renewcommand{\cite}{\textcite}
\newcommand{\citet}{\textcite}

\newcommand{\np}{\newpage}
\newcommand{\rt}[1]{Theorem~\ref{#1}}
\newcommand{\bpf}{\begin{proof}}
\newcommand{\epf}{\end{proof}}
\newcommand{\nn}{\nonumber}
\newcommand{\beq}{\begin{equation}}
\newcommand{\eeq}{\end{equation}}
\newcommand{\beqn}{\[}
\newcommand{\eeqn}{\]}
%\newcommand{\beq}{\begin{eqnarray}}
%\newcommand{\eeq}{\end{eqnarray}}
%\newcommand{\beqn}{\begin{eqnarray*}}
%\newcommand{\eeqn}{\end{eqnarray*}}
\newcommand{\reqn}[1]{Equation~(\ref{#1})}
\newcommand{\reprop}[1]{Proposition~\ref{#1}}
\newcommand{\bit}{\begin{itemize}}
\newcommand{\eit}{\end{itemize}}
\newcommand{\ben}{\begin{enumerate}}
\newcommand{\een}{\end{enumerate}}
\newcommand{\ignore}[1]{{}}
\newcommand{\rem}{Remark:}
\newcommand{\rfig}[1]{{Figure~\ref{#1}}}
\newcommand{\rtab}[1]{{Table~\ref{#1}}}
\newcommand{\rqt}[1]{Q.~\ref{#1}}
\renewcommand{\it}{\item}

\newtheorem{defn}{Definition}[chapter]
\newtheorem{thm}{Theorem}[chapter]
\newtheorem{fact}{Fact}[chapter]
\newtheorem{lem}{Lemma}
\newtheorem{prop}{Proposition}
\newtheorem{cor}{Corollary}
\newtheorem{note}{Note}[chapter]
\newtheorem{claim}{Claim}
\newtheorem{step}{Step}
\newtheorem{ex}{{\em Example}}[chapter]
\newtheorem{assumption}{Assumption}
\newtheorem{approximation}{Approximation}
\newtheorem{property}{Property}[chapter]
\newtheorem{obs}{Observation}[chapter]
\newtheorem{exer}{{\em Exercise}}[chapter]

%% Electromagnetics
\newcommand{\del}{\nabla}
\newcommand{\pd}[2]{\frac{\partial{#1}_{#2}}{\partial #2}}
\newcommand{\pda}[2]{\frac{\partial #1}{\partial #2}}
\newcommand{\pdda}[2]{\frac{\partial^2 #1}{\partial {#2}^2}}
\newcommand{\dif}[2]{\frac{d #1}{d #2}}
\newcommand{\ddif}[2]{\frac{d^2 #1}{d #2^2}}
\newcommand{\sm}[1]{\triangle #1}

\newcommand{\uvec}[1]{\vec{\mathbf{a}}_{#1}}
%\newcommand{\uvec}[1]{\mathbf{a}_{#1}}
\newcommand{\cvec}[1]{\vec{\mathbf{#1}}}
\newcommand{\diff}{\mathrm{d}}

\newcommand{\ux}{\uvec{x}}
\newcommand{\uy}{\uvec{y}}
\newcommand{\uz}{\uvec{z}}
\newcommand{\ur}{\uvec{r}}
\newcommand{\uR}{\uvec{R}}
\newcommand{\utheta}{\uvec{\theta}}
\newcommand{\uphi}{\uvec{\phi}}
\newcommand{\urho}{\uvec{\rho}}
\newcommand{\un}{\uvec{n}}

\newcommand{\va}{\cvec{A}}
\newcommand{\vA}{\cvec{A}}
\newcommand{\vb}{\cvec{B}}
\newcommand{\vB}{\cvec{B}}
\newcommand{\vdb}{\cvec{d\mathbf{B}}}
\newcommand{\vc}{\cvec{C}}
\newcommand{\ve}{\cvec{E}}

%\newcommand{\myframe}[1]{\begin{framed} #1 \end{framed}}
\newcommand{\mf}[1]{\begin{framed} #1 \end{framed}}
\newcommand{\hili}[1]{\textbf{\color{red} #1}}
\newcommand{\bex}[1]{\mf{\begin{ex} #1 \end{ex}}}
\newcommand{\bexer}[1]{\mf{\begin{exer} #1 \end{exer}}}
\newcommand{\bnote}[1]{\mf{\begin{note} #1 \end{note}}}

\newcommand{\bfig}[3]{\begin{figure}[H] \centering \includegraphics[width=#1\textwidth]{#2} 	\caption{#3} \end{figure}}

\newcommand{\bfigl}[3]{\begin{figure}[H] \centering \includegraphics[width=#1\textwidth]{#2} 	\caption{#3} \label{#2} \end{figure}}

\newcommand{\bfiglrotate}[4]{\begin{figure}[H] \centering \includegraphics[angle=#4,origin=c,width=#1\linewidth]{#2} 	\caption{#3} \label{#2} \end{figure}}

\newcommand{\bfigll}[4]{\begin{figure}[H] \includegraphics[width=#1\textwidth]{#2} 	\caption{#3} \label{#4} \end{figure}}

\newcommand{\bfigtwo}[5]{
\begin{figure}[H]
\begin{subfigure}{#1\textwidth}
  \centering  \includegraphics[width=1\linewidth]{#2} 
  \caption{} 
\end{subfigure} \hfill
\begin{subfigure}{#3\textwidth}
  \centering  \includegraphics[width=1\linewidth]{#4} 
  \caption{} 
\end{subfigure}\hfill
\caption{#5}
\label{#2}
\end{figure}
}

% \makeatletter
% \newcommand{\unchapter}[1]{% FROM https://tex.stackexchange.com/a/624410/161015
%    \begingroup
%    \let\@makeschapterhead\@gobble % make \@makechapterhead do nothing
%    \chapter*{#1}
%    \addcontentsline{toc}{chapter}{#1}
%    \markboth{#1}{}
%    \endgroup
%}
%\makeatother